% !TEX root = ../main.tex
\chapter{Codimensión, Dicotomía topológica e índice \texorpdfstring{$Q_k$}{Qk}}\label{sec:codimension_qk}

\section{Codimensión topológica}

Partimos del producto interno de Sobolev introducido en el Capítulo~\ref{sec:intro}. Denotamos por $\mathcal H$ la completación de $\Poly$ respecto de la norma inducida por ese producto. En particular,
\[
    \mathcal H=\overline{\Poly},
\]
y $\Poly$ es un subespacio denso de $\mathcal H$.

La noción clásica de codimensión en este contexto es la siguiente.

\begin{defn}[Codimensión topológica en Hilbert \cite{halmos1982hilbert}]\label{def:codim_top_hilbert}
    Sea $\mathcal H$ un espacio de Hilbert y sea $M\subset \mathcal H$ un subespacio cerrado. Definimos
    \[
        \operatorname{codim}_{\mathcal H}(M):=\dim(\mathcal H/M).
    \]
    Equivalentemente,
    \[
        \operatorname{codim}_{\mathcal H}(M)=\dim(M^\perp).
    \]
\end{defn}

\begin{prop}[Dicotomía topológica para un funcional]\label{thm:dicotomia_1}
    Sea $\phi:\Poly\to\mathbb C$ un funcional lineal no nulo y $(\mathcal{H}, \langle\cdot, \cdot\rangle)$. Entonces ocurre exactamente una de las dos posibilidades siguientes:
    \begin{enumerate}
        \item si $\phi$ no es acotado respecto de la norma de $\mathcal H$, entonces
              \[
                  \overline{\ker\phi}=\mathcal H,
                  \qquad
                  \operatorname{codim}_{\mathcal H}(\overline{\ker\phi})=0;
              \]
        \item si $\phi$ es acotado respecto de la norma de $\mathcal H$, entonces
              \[
                  \operatorname{codim}_{\mathcal H}(\overline{\ker\phi})=1.
              \]
    \end{enumerate}
\end{prop}

\begin{proof}
    Si $\phi$ no es acotado, su núcleo es denso en $\Poly$ respecto de la norma inducida por $\mathcal H$. En efecto, fijado $p\in\Poly$, como $\phi$ es no acotado existe una sucesión $u_n\in\Poly$ con $\|u_n\|\to 0$ y $\phi(u_n)=1$. Entonces
    \[
        v_n:=p-\phi(p)u_n\in\ker\phi,
        \qquad
        \|v_n-p\|=|\phi(p)|\,\|u_n\|\to 0.
    \]
    Como $\Poly$ es denso en $\mathcal H$, se sigue que $\overline{\ker\phi}=\mathcal H$.

    Supongamos ahora que $\phi$ es acotado. Entonces se extiende de manera única a un funcional continuo, también denotado por $\phi$, sobre $\mathcal H$. Su núcleo es cerrado y contiene a $\ker\phi\subset\Poly$, luego
    \[
        \overline{\ker\phi}^{\,\mathcal H}=\ker\phi.
    \]
    Tomemos $q\in\Poly$ con $\phi(q)\neq 0$. Para todo $x\in\mathcal H$,
    \[
        x=\left(x-\frac{\phi(x)}{\phi(q)}q\right)+\frac{\phi(x)}{\phi(q)}q,
    \]
    y el primer sumando pertenece a $\ker\phi$. Por tanto,
    \[
        \mathcal H=\ker\phi+\operatorname{span}\{q\}.
    \]
    Como $q\notin\ker\phi$, esa suma es directa. En consecuencia,
    \[
        \operatorname{codim}_{\mathcal H}(\overline{\ker\phi})=\operatorname{codim}_{\mathcal H}(\ker\phi)=1.
    \]
\end{proof}

\begin{prop}
    Sea $\mathcal{H}$ un espacio de Hilbert y sean $\phi_1, \dots, \phi_N$ funcionales lineales continuos y linealmente independientes. Entonces, el subespacio cerrado
    \[
        V = \bigcap_{i=1}^N \ker(\phi_i)
    \]
    tiene codimensión topológica $\operatorname{codim}_{\mathcal{H}}(V) = N$.
\end{prop}

\begin{proof}
    Al ser cada $\phi_i$ continuo en $\mathcal{H}$, por el Teorema de Riesz existen representantes $f_1, \dots, f_N \in \mathcal{H}$ linealmente independientes tales que $\phi_i(x) = \langle x, f_i \rangle_{\mathcal{H}}$. El subespacio $V$ es el complemento ortogonal del subespacio $F = \operatorname{span}\{f_1, \dots, f_N\}$. Como $\dim(F) = N$, se concluye que $\operatorname{codim}_{\mathcal{H}}(V) = N$.
\end{proof}

\section{Codimensión algebraica}

Dado que la codimensión topológica puede colapsar al clausurar subespacios definidos por funcionales no continuos, necesitamos una noción que preserve exactamente el número de condiciones lineales impuestas.

\begin{defn}\label{def:codim_alg}
    Sea $V\subset \Poly$ un subespacio vectorial. Definimos
    \[
        \operatorname{codim}_{\Poly}(V):=\dim(\Poly/V)\in\mathbb N\cup\{\infty\},
    \]
    donde la dimensión es algebraica, es decir, la cardinalidad de una base de Hamel \cite{hormander1989linearfunctionalanalysis,Aizpuru2009Apuntes}.
\end{defn}

La primera caracterización es la que usaremos de manera sistemática.

\begin{prop}\label{prop:finite-codim-kernels}
    Para un subespacio vectorial $V\subset \Poly$, son equivalentes:
    \begin{enumerate}
        \item $\operatorname{codim}_{\Poly}(V)<\infty$;
        \item existe un sistema de $m\in\mathbb N$ funcionales lineales $\{\phi_1,\dots,\phi_m\}$ tales que
              \[
                  V=\bigcap_{j=1}^m \ker \phi_j.
              \]
    \end{enumerate}
    Además,
    \[
        \operatorname{codim}_{\Poly}(V)
        =
        \min\Bigl\{m:\exists\,\{\phi_1,\dots,\phi_m\},\ V=\bigcap_{j=1}^m\ker\phi_j\Bigr\},
    \]
    y más precisamente, si $V=\bigcap_{j=1}^m\ker\phi_j$, entonces
    \[
        \operatorname{codim}_{\Poly}(V)=\dim\operatorname{span}\{\phi_1,\dots,\phi_m\}.
    \]
\end{prop}

\begin{proof}
    Probemos primero que $(1)\Rightarrow(2)$. Si
    \(
    \operatorname{codim}_{\Poly}(V)=n<\infty,
    \)
    existe un subespacio complementario $F\subset\Poly$ tal que
    \[
        \Poly=V\oplus F,
        \qquad
        \dim(F)=n.
    \]
    Tomemos una base $u_1,\dots,u_n$ de $F$ y sea $\lambda_1,\dots,\lambda_n$ la base dual de $F^*$. Para cada $p\in\Poly$, escribiendo de manera única $p=v+f$ con $v\in V$ y $f\in F$, definimos
    \[
        \phi_j(p):=\lambda_j(f),\qquad j=1,\dots,n.
    \]
    Entonces cada $\phi_j$ es lineal y
    \[
        V=\bigcap_{j=1}^n\ker\phi_j,
    \]
    pues $\phi_j(p)=0$ para todo $j$ si y solo si la componente $f$ es nula.

    Probemos ahora $(2)\Rightarrow(1)$. Supongamos que
    \[
        V=\bigcap_{j=1}^m\ker\phi_j.
    \]
    Sea
    \[
        r:=\dim\operatorname{span}\{\phi_1,\dots,\phi_m\}.
    \]
    Reordenando si hace falta, podemos suponer que $\phi_1,\dots,\phi_r$ forman una base de ese subespacio. Entonces
    \[
        V=\bigcap_{i=1}^r\ker\phi_i.
    \]
    Consideremos la aplicación lineal
    \[
        \Psi:\Poly\to\mathbb C^r,
        \qquad
        \Psi(p)=(\phi_1(p),\dots,\phi_r(p)).
    \]
    Como $\phi_1,\dots,\phi_r$ son linealmente independientes, $\Psi$ es sobreyectiva. Además,
    \[
        \ker\Psi=V.
    \]
    Por el primer teorema de isomorfía,
    \[
        \Poly/V\cong\mathbb C^r,
    \]
    y por tanto
    \[
        \operatorname{codim}_{\Poly}(V)=r=
        \dim\operatorname{span}\{\phi_1,\dots,\phi_m\}.
    \]
    La fórmula del mínimo se deduce de inmediato de esta identidad.
\end{proof}

A continuación damos una caracterización de la codimensión algebraica atendiendo a la ortogonalidad dada por el producto escalar.

\begin{prop}\label{prop:equiv_ortho_interseccion}
    Sea $\{f_1,\dots,f_m\}\subset \Poly$ y $F=\operatorname{span}\{f_1,\dots,f_m\}$ y supongamos que $\langle\cdot,\cdot\rangle$ es un producto escalar en $\Poly$. Definimos
    \[
        F^{\perp}:=\{p\in\Poly:\langle p,f_j\rangle=0,\ j=1,\dots,m\}.
    \]
    Entonces
    \[
        F^{\perp}=\bigcap_{j=1}^m \ker(\phi_j),
        \qquad
        \phi_j(p):=\langle p,f_j\rangle,
    \]
    y
    \[
        \operatorname{codim}_{\Poly}(F^{\perp})
        =
        \dim F.
    \]
    En particular, si $f_1,\dots,f_m$ son linealmente independientes, entonces
    \[
        \operatorname{codim}_{\Poly}(F^{\perp})=m.
    \]
\end{prop}

\begin{proof}
    La identidad
    \[
        F^{\perp}=\bigcap_{j=1}^m\ker(\phi_j)
    \]
    es inmediata de la definición.

    Sea ahora
    \[
        r:=\dim\operatorname{span}\{f_1,\dots,f_m\}.
    \]
    Reordenando si hace falta, podemos suponer que $f_1,\dots,f_r$ forman una base de ese subespacio. Entonces, para cada $j>r$, el vector $f_j$ es combinación lineal de $f_1,\dots,f_r$, y por sesquilinealidad del producto interno se sigue que $\phi_j$ es combinación lineal de $\phi_1,\dots,\phi_r$. Por tanto,
    \[
        F^{\perp}=\bigcap_{i=1}^r\ker(\phi_i).
    \]

    Además, los funcionales $\phi_1,\dots,\phi_r$ son linealmente independientes. En efecto, si
    \[
        \sum_{i=1}^r \alpha_i\phi_i=0,
    \]
    entonces, para todo $p\in\Poly$,
    \[
        0=\sum_{i=1}^r\alpha_i\langle p,f_i\rangle
        =\Bigl\langle p,\sum_{i=1}^r\overline{\alpha_i}f_i\Bigr\rangle.
    \]
    Tomando
    \(
    p=\sum_{i=1}^r\overline{\alpha_i}f_i
    \)
    obtenemos
    \[
        \Bigl\|\sum_{i=1}^r\overline{\alpha_i}f_i\Bigr\|^2=0,
    \]
    luego $\alpha_1=\cdots=\alpha_r=0$.

    Aplicando la Proposición~\ref{prop:finite-codim-kernels} a $\phi_1,\dots,\phi_r$, concluimos que
    \[
        \operatorname{codim}_{\Poly}(V_\perp(F))=r=
        \dim\operatorname{span}\{f_1,\dots,f_m\}.
    \]
\end{proof}

En este trabajo son relevantes los siguientes funcionales de evaluación y evaluación de la derivada en un punto.

\begin{defn}\label{def:funcionales_restringidos}
    Sea $c\in\C$. Definimos los funcionales lineales
    \[
        \delta_c:\Poly\to\C,
        \qquad
        \delta_c(p)=p(c),
    \]
    y
    \[
        \delta'_c:\Poly\to\C,
        \qquad
        \delta'_c(p)=p'(c).
    \]
\end{defn}

\begin{lem}
    Sea $\mathcal{A} = \{z_1, \dots, z_N\} \subset \C$ un conjunto finito de $N$ puntos distintos. Entonces el conjunto de funcionales de evaluación $\{\delta_{z_1}, \dots, \delta_{z_N}\}$ es linealmente independiente en el espacio dual de funcionales lineales$\Poly^*$.
\end{lem}

\begin{proof}
    Supongamos que existe una combinación lineal nula $\sum_{i=1}^N \alpha_i \delta_{z_i} = 0$.
    Para cada índice $j \in \{1, \dots, N\}$, construimos el polinomio de interpolación de Lagrange:
    \[
        L_j(z) = \prod_{\substack{l=1 \\ l \neq j}}^{N} \frac{z - z_l}{z_j - z_l}
    \]
    Satisface que $L_j(z_i) = \delta_{ji}$ (Delta de Kronecker). Evaluando la combinación lineal nula en el polinomio $L_j$, obtenemos:
    \[
        0 = \sum_{i=1}^N \alpha_i \delta_{z_i}(L_j) = \sum_{i=1}^N \alpha_i L_j(z_i) = \alpha_j \cdot 1
    \]
    Por tanto, $\alpha_j = 0$ para todo $j$, lo que demuestra la independencia.
\end{proof}

\begin{cor}
    Sean $z_1, \dots, z_N \in \C$ puntos todos distintos. El subespacio cerrado de los polinomios que se anulan en dichos puntos, dado por
    \[
        V = \bigcap_{i=1}^N \ker(\delta_{z_i}),
    \]
    tiene codimensión algebraica exacta $N$ en $\Poly$.
\end{cor}

\begin{proof}
    Por el lema anterior, el sistema de funcionales de evaluación $\{\delta_{z_1}, \dots, \delta_{z_N}\}$ es linealmente independiente. Como $V$ está definido como la intersección de los núcleos de estos $N$ funcionales, la Proposición~\ref{prop:finite-codim-kernels} establece que su codimensión algebraica es exactamente la dimensión del subespacio que generan. Por tanto,
    \[
        \operatorname{codim}_{\Poly}(V) = \dim\operatorname{span}\{\delta_{z_1}, \dots, \delta_{z_N}\} = N.
    \]
\end{proof}

\section{Relación entre codimensión algebraica y topológica}

En esta sección estudiaremos la relación entre las codimensiones algebraicas y topológicas introducidas en las secciones anteriores.

\begin{prop}\label{prop:relacion_codimensiones}
    Sea $\mathcal H=\overline{\Poly}$ con respecto a un producto interno.
    \begin{enumerate}
        \item Si $V\subset\Poly$ satisface $\operatorname{codim}_{\Poly}(V)=m<\infty$, entonces
              \[
                  \operatorname{codim}_{\mathcal H}(\overline V)\le m.
              \]
        \item Si $M\subset\mathcal H$ es cerrado y $\operatorname{codim}_{\mathcal H}(M)=m<\infty$, entonces
              \[
                  \overline{M\cap\Poly}=M
                  \qquad\text{y}\qquad
                  \operatorname{codim}_{\Poly}(M\cap\Poly)=m.
              \]
    \end{enumerate}
\end{prop}

\begin{proof}
    Probemos $(1)$. Tomemos un subespacio complementario $F\subset\Poly$ con
    \[
        \Poly=V\oplus F,
        \qquad
        \dim(F)=m.
    \]
    Como $F$ es finito-dimensional, es cerrado en $\mathcal H$. Al tomar clausuras,
    \[
        \mathcal H=\overline{\Poly}=\overline{V\oplus F}=\overline V+F.
    \]
    La aplicación cociente $F\to \mathcal H/\overline V$, $f\mapsto f+\overline V$, es sobreyectiva. Luego
    \[
        \dim(\mathcal H/\overline V)\le \dim(F)=m.
    \]

    Probemos $(2)$. Como $M$ es cerrado y de codimensión $m$, existen $m$funcionales lineales continuos $\Phi_1,\dots,\Phi_m$ tal que
    \[
        M=\bigcap_{j=1}^m \ker \Phi_j,
    \]
    Restringiendo a $\Poly$, obtenemos
    \[
        M\cap\Poly=\bigcap_{j=1}^m \ker(\Phi_j|_{\Poly}).
    \]
    Las restricciones $\Phi_j|_{\Poly}$ son linealmente independientes: si una combinación lineal se anula en $\Poly$, también se anula en todo $\mathcal H$ por densidad, puesto que todos los funcionales son continuos; y entonces todos los coeficientes son nulos. Por la Proposición~\ref{prop:finite-codim-kernels},
    \[
        \operatorname{codim}_{\Poly}(M\cap\Poly)=m.
    \]
    Además, $M\cap\Poly$ es denso en $M$. En efecto, fijado $x\in M$, tomamos una sucesión $p_n\in\Poly$ tal que $p_n\to x$. Definimos
    \[
        a_n:=(\Phi_1(p_n),\dots,\Phi_m(p_n))\in\mathbb C^m.
    \]
    Como $x\in M$, se tiene $a_n\to 0$, ya que al ser los funcionales continuos, $\Phi_j(p_n)\to\Phi_j(x)=0$ para cada $j$. Consideremos ahora la aplicación lineal
    \[
        S:\Poly\to\mathbb C^m,
        \qquad
        S(p)=(\Phi_1(p),\dots,\Phi_m(p)).
    \]
    Las restricciones $\Phi_j|_{\Poly}$ son linealmente independientes, luego $S$ es sobreyectiva. Elegimos polinomios $u_1,\dots,u_m\in\Poly$ tales que $S(u_i)=e_i$, donde $e_i$ es la base canónica de $\mathbb C^m$, y definimos
    \[
        R:\mathbb C^m\to\Poly,
        \qquad
        R(b_1,\dots,b_m):=\sum_{i=1}^m b_i u_i.
    \]
    Entonces $S\circ R=\mathrm{Id}_{\mathbb C^m}$. Para cada $n$, ponemos
    \[
        v_n:=p_n-R(a_n)\in\Poly.
    \]
    Como $S(v_n)=0$, se sigue que $v_n\in M\cap\Poly$. Además,
    \[
        \|v_n-x\|\le \|p_n-x\|+\|R(a_n)\|\xrightarrow[n\to\infty]{}0,
    \]
    porque $p_n\to x$, $a_n\to0$ y $R$ es continua sobre el espacio finito-dimensional $\mathbb C^m$. Luego $x\in\overline{M\cap\Poly}$, y por tanto $\overline{M\cap\Poly}=M$.
\end{proof}

\begin{rmk}
    La desigualdad de la parte $(1)$ puede ser estricta. Eso es exactamente lo que muestra la Proposición~\ref{thm:dicotomia_1}: un subespacio de codimensión algebraica uno puede tener clausura total en $\mathcal H$ y, por tanto, codimensión topológica de su adherencia es $0$.
\end{rmk}

\section{Índice \texorpdfstring{$Q_k$}{Qk}}

Recordemos que dado un operador lineal $T \in \mathcal{L}(\mathcal{H})$ un subespacio vectorial $M$ se define la restricción de un operador $T$ como $T|_M$ y la norma es $\|T|_M\|=\sup_{\substack{p\in M\\ p\ne0}}\frac{\|Tp\|}{\|p\|} \leq \|T\| < +\infty$

\begin{defn}\label{def:numeros_gelfand}
    Sea $T\in\mathcal L(\mathcal H)$. Para cada $k\ge 0$, definimos
    \begin{equation}\label{eq:gelfand_numbers}
        c_k(T)=\inf\{\|T|_M\|:M\subset \mathcal H,\ M\text{ cerrado},\ \operatorname{codim}_{\mathcal H}(M)<k+1\}.
    \end{equation}
    En particular, $c_0(T)=\|T\|$. \cite{Ullrich2024}
\end{defn}
\begin{rmk}
    Nuestro índice $c_0(T)$ corresponde con el índice $c_1(T)$, y en general $c_{k-1}(T)$ será $c_k(T)$. Lo definimos de esta manera para coherencia en notación al trabajar en espacios de polinomios.
\end{rmk}

El índice de Gelfand se define para operadores acotados, en este trabajo aparecerán operadores no acotados. Esto justifica la introducción de un nuevo índice para lo que tenemos que introducir una nueva noción de codimensión, de la que hablaremos más adelante.

A continuación introducimos un nuevo índice que extiende la noción de los números de Gelfand a operadores no acotados sobre $\Poly$.
\begin{defn}\label{def:indice_Qk}
    Sea $T:\Poly\to\Poly$ un operador lineal no necesariamente acotado definido sobre $(\Poly,\|\cdot\|)$, donde la norma está inducida por un producto interno. Para cada $k\ge0$, definimos
    \[
        Q_k(T)
        :=\inf\bigl\{\|T|_V\|:V\subset\Poly\ \ \operatorname{codim}_{\Poly}(V)<k+1\bigr\},
    \]
    donde $Q_k(T)\in [0,+\infty]$.
\end{defn}

\begin{rmk}
    En particular, para $k=0$ la condición
    \(
    \operatorname{codim}_{\Poly}(V)<1
    \)
    fuerza $V=\Poly$, y por tanto
    \[
        Q_0(T)=\|T\|,
    \]
    admitiendo el valor $+\infty$ cuando $T$ no es acotado.
\end{rmk}

\begin{prop}\label{prop:propiedades_Qk}
    Sea $T:\Poly\to\Poly$ un operador lineal. Entonces
    \[
        \|T\|=Q_0(T)\ge Q_1(T)\ge \cdots \ge Q_k(T)\ge Q_{k+1}(T)\ge \cdots \ge 0,
    \]
    admitiendo el valor $+\infty$ cuando corresponda.
\end{prop}

\begin{proof}
    La identidad $Q_0(T)=\|T\|$ es inmediata de la definición.

    Para probar la monotonía, consideremos la familia de subespacios admisibles para el índice $k$:
    \[
        \mathcal{F}_k = \{ V \subset \Poly : \operatorname{codim}_{\Poly}(V) < k+1 \}.
    \]
    Si un subespacio $V$ pertenece a $\mathcal{F}_k$, su codimensión es como máximo $k$, lo que implica que también es estrictamente menor que $k+2$. Por tanto, el subespacio también es admisible para $k+1$, de modo que tenemos la inclusión de conjuntos $\mathcal{F}_k \subset \mathcal{F}_{k+1}$.

    Como $Q_k(T) = \inf_{V \in \mathcal{F}_k} \|T|_V\|$, tomando el ínfimo:
    \[
        Q_{k+1}(T) \le Q_k(T).
    \]
    La no negatividad es trivial.
\end{proof}

\begin{defn}
    Sea $T:\Poly\to\Poly$ un operador lineal no necesariamente acotado, diremos que la cadena 
    \[
        Q_0(T)\ge Q_1(T)\ge \cdots \ge Q_k(T)\ge Q_{k+1}(T)\ge \cdots \ge 0,
    \]
    estabiliza en $k$ si $Q_{k-1}(T)=\infty$ y $Q_k(T)<\infty$.
\end{defn}

A continuación fijamos un modelo elemental en el que la cadena se rompe por primera vez en $k=1$.

\begin{exm}
    Consideremos el producto de Sobolev
    \[
        \langle p,q\rangle_S=\int_{\mathbb S_1}p(z)\overline{q(z)}\,d\mathbf m(z)+p'(2)\overline{q'(2)},
        \qquad p,q\in\Poly,
    \]
    y el operador de multiplicación $\mathcal D p(z)=zp(z)$. Para analizar su acotación, calculamos la norma de $\mathcal D p$:
    \[
        \|\mathcal D p\|_S^2=\int_{\mathbb S_1}|zp(z)|^2\,d\mathbf m(z)+|(\mathcal D p)'(2)|^2=\|p\|_{L^2(\mathbf m)}^2+|p(2)+2p'(2)|^2.
    \]
    Podemos construir una sucesión explícita de polinomios $\{p_n\}_{n\in\mathbb{N}}$ para comprobar que $\mathcal D$ no es acotado. Consideremos:
    \[
        p_n(z) = \frac{n^2}{2^n} \left( z^n - \frac{n}{2(n+1)} z^{n+1} \right).
    \]
    Es directo comprobar que su derivada se anula en $z=2$:
    \[
        p_n'(z) = \frac{n^2}{2^n} \left( n z^{n-1} - \frac{n}{2} z^n \right) \implies p_n'(2) = \frac{n^2}{2^n} \left( n 2^{n-1} - n 2^{n-1} \right) = 0.
    \]
    Recordando que los monomios son ortonormales en $L^2(\mathbf m)$, su norma es:
    \[
        \|p_n\|_{L^2(\mathbf m)}^2 = \frac{n^4}{4^n} \left( 1 + \frac{n^2}{4(n+1)^2} \right) \xrightarrow[n\to\infty]{} 0.
    \]
    Sin embargo, su evaluación en $z=2$ diverge:
    \[
        p_n(2) = \frac{n^2}{2^n} \left( 2^n - \frac{n}{2(n+1)} 2^{n+1} \right) = n^2 \left( 1 - \frac{n}{n+1} \right) = \frac{n^2}{n+1} \xrightarrow[n\to\infty]{} \infty.
    \]
    Para esta sucesión se tiene que
    \[
        \|p_n\|_S^2 = \|p_n\|_{L^2(\mathbf m)}^2 + |p_n'(2)|^2 = \|p_n\|_{L^2(\mathbf m)}^2 \to 0,
    \]
    mientras que
    \[
        \|\mathcal D p_n\|_S^2 = \|p_n\|_{L^2(\mathbf m)}^2 + |p_n(2)|^2 \to \infty.
    \]
    Esto demuestra que el operador $\mathcal D$ no es acotado globalmente, es decir, $Q_0(\mathcal D)=\infty$. 
    
    Sin embargo, la no acotación proviene del término libre $p(2)$. Si consideramos el subespacio de codimensión algebraica 1:
    \[
        V = \{p \in \Poly : p(2) = 0\},
    \]
    para cualquier $p \in V$ se tiene que $(\mathcal D p)'(2) = p(2) + 2p'(2) = 2p'(2)$. Evaluando nuevamente la norma obtenemos:
    \[
        \|\mathcal D p\|_S^2 = \|p\|_{L^2(\mathbf m)}^2 + |2p'(2)|^2 \le 4\left(\|p\|_{L^2(\mathbf m)}^2 + |p'(2)|^2\right) = 4\|p\|_S^2.
    \]
    Por tanto, la restricción del operador a este subespacio es acotada, lo que implica que
    \[
        Q_1(\mathcal D) \le 2 < \infty.
    \]
\end{exm}

De este ejemplo es importante observar que existen operadores lineales no acotados sobre $\Poly$ y sin embargo su restricción a subespacios ``grandes'' (en el sentido de que tienen codimensión 1)  es acotado. Esto justifica el estudio de la estabilización de la cadena de los $Q_k$ para encontrar subespacios de codimensión finita donde los operadores son acotados.

\section{Equivalencia con los números de Gelfand}

A continuación, probaremos que cuando consideramos operadores lineales acotados sobre $\Poly$, que admiten por tanto una extensión al espacio de Hilbert $\mathcal H$, entonces los índices $Q_k$ coinciden con los números de Gelfand.

\begin{thm}\label{thm:equivalencia_Qk_Gelfand}
    Sea $(\mathcal H,\langle\cdot,\cdot\rangle)$ un espacio de Hilbert tal que $\mathcal H=\overline{\Poly}$, y sea $T\in\mathcal L(\mathcal H)$ con $T(\Poly)\subset\Poly$. Entonces, para todo $k\ge0$,
    \begin{equation}\label{eq:igualdad_Qk_Gelfand}
        Q_k(T)=c_k(T),
    \end{equation}
    donde $Q_k(T)$ se calcula sobre $\Poly$ con la norma inducida por $\mathcal H$.
\end{thm}

\begin{proof}
    Probemos primero que $c_k(T)\le Q_k(T)$. Sea $V\subset\Poly$ tal que
    \[
        \operatorname{codim}_{\Poly}(V)=m<k+1.
    \]
    Poniendo $M:=\overline V$, la Proposición~\ref{prop:relacion_codimensiones} da
    \[
        \operatorname{codim}_{\mathcal H}(M)\le m<k+1.
    \]
    Por tanto, $M$ es un subespacio admisible en la definición de $c_k(T)$. Como $T$ es continuo y $V$ es denso en $M$,
    \[
        \|T|_M\|=\|T|_V\|.
    \]
    Luego
    \[
        c_k(T)\le \|T|_M\|=\|T|_V\|.
    \]
    Tomando ínfimo sobre todos esos $V$, obtenemos
    \[
        c_k(T)\le Q_k(T).
    \]

    Probemos ahora que $Q_k(T)\le c_k(T)$. Sea $M\subset\mathcal H$ cerrado con
    \[
        \operatorname{codim}_{\mathcal H}(M)=m<k+1.
    \]
    Entonces, por la Proposición~\ref{prop:relacion_codimensiones}, el subespacio
    \[
        V:=M\cap\Poly
    \]
    satisface
    \[
        \operatorname{codim}_{\Poly}(V)=m,
        \qquad
        \overline V=M,
    \]
    y nuevamente
    \[
        \|T|_V\|=\|T|_M\|.
    \]
    Por tanto,
    \[
        Q_k(T)\le \|T|_V\|=\|T|_M\|.
    \]
    Tomando ínfimo sobre todos los subespacios $M$ admisibles en la definición de $c_k(T)$, llegamos a
    \[
        Q_k(T)\le c_k(T).
    \]

    Las dos desigualdades prueban \eqref{eq:igualdad_Qk_Gelfand}.
\end{proof}

\begin{rmk}
    Esto justifica que el índice $Q_k$ es una extensión del índice de Gelfand para operadores no acotados, ya que coinciden cuando estos son acotados.
\end{rmk}
